\documentclass[10pt]{article}

\usepackage[margin=0.9in]{geometry}
\usepackage{enumitem}
\usepackage{booktabs}

\setlength{\parskip}{0.25em}
\setlist{nosep, leftmargin=*}

\title{\vspace{-1.5em}ECE 268 Project Progress Report\\[0.2em]
\large Lightweight Block Ciphers for Constrained Devices: PRESENT-80 and SPECK-64}
\author{%
  Devansh Gupta \and
  Priyansh Parikh \and
  Jeevan N V
}
\date{May 21, 2026}

\begin{document}
\maketitle
\vspace{-1.2em}

\section*{Project Summary}
\textbf{Goal:} Implement PRESENT-80 and SPECK-64 from scratch with encryption, decryption,
and key schedules. Each cipher will be wrapped in CBC and CTR modes, validated against
published test vectors, and benchmarked on GPU against a sequential CPU baseline and
AES-128, with IoT/RFID as the target setting where AES is too expensive. \\
\textbf{Status:} PRESENT-80 is complete on both CPU and GPU. The GPU implementation
is written in pure CUDA C++ compiled directly with \texttt{nvcc}, rather than going
through PyCUDA or CuPy wrappers. All four official test vectors pass and GPU
performance has been measured. SPECK-64, CBC/CTR modes, and the AES baseline are
pending.

\section{Work Completed So Far}

\subsection*{Implementations}
\begin{itemize}
  \item \textbf{Repository:} Current files are \texttt{present80.cpp}, \texttt{present80.cu},
        and \texttt{CMakeLists.txt}. The planned files \texttt{speck64.cpp},
        \texttt{speck64.cu}, \texttt{modes.cpp}, and an AES baseline are not yet written.

  \item \textbf{PRESENT-80 CPU:} 64-bit block, 80-bit key, 32 rounds per Bogdanov et al.\
        The key schedule rotates the 80-bit register left by 61, applies the S-box to
        the top nibble, and XORs the round counter into bits 19--15. Encryption runs 31
        full rounds of addRoundKey, sBoxLayer, and pLayer, then adds the final round key.
        Decryption reverses those steps using INV\_SBOX and the inverse P-layer. The bit
        permutation formula is $p(j) = (16j) \bmod 63$ with bit 63 fixed in place.

  \item \textbf{PRESENT-80 GPU:} A pure-CUDA C++ port of the CPU implementation,
        written directly against the CUDA Runtime API and built with \texttt{nvcc}
        --- no PyCUDA, CuPy, or other higher-level wrapper sits between the host
        code and the device. Round keys and S-box tables go into
        \texttt{\_\_constant\_\_} memory so they are broadcast across each warp
        without extra traffic. The encrypt and decrypt kernels each assign one
        thread per cipher block. The P-layer helper is
        \texttt{\_\_device\_\_ \_\_forceinline\_\_} with \texttt{\#pragma unroll}
        on the 63-iteration loop. Key expansion runs on the CPU and the result is
        copied to constant memory before the kernel launches. The test bench
        encrypts and then decrypts $2^{20}$ distinct blocks and verifies that every
        output matches the original plaintext.

  \item \textbf{Build system:} \texttt{CMakeLists.txt} uses a \texttt{USE\_CUDA} flag.
        \texttt{ON} builds \texttt{present80.cu} with \texttt{nvcc} targeting the native
        GPU architecture. \texttt{OFF} builds \texttt{present80.cpp} as a C++17 binary
        with \texttt{-O4}.
\end{itemize}

\subsection*{Experiments and Evaluations}
\begin{itemize}
  \item \textbf{Setup:} Test vectors are the four canonical cases from Bogdanov et al.\
        GPU timing uses CUDA events over $2^{20}$ blocks on an NVIDIA RTX~4070 Mobile
        running CUDA~12.0. CPU timing uses \texttt{std::chrono::high\_resolution\_clock}
        on the same machine at $\approx$3.8~GHz.

  \item \textbf{Completed:} All four PRESENT-80 spec vectors pass on the CPU build.
        GPU round-trip correctness was verified across all $2^{20}$ blocks. CPU and GPU
        encryption throughput were both measured on the same workload.

  \item \textbf{Pending:} SPECK-64 and CBC/CTR tests have not been run because those
        components are not yet implemented.
\end{itemize}

\subsection*{Preliminary Results}

\begin{table}[h]
\centering
\small
\begin{tabular}{lcccc}
\toprule
\textbf{Cipher / Platform} & \textbf{Block/Key} & \textbf{Enc/Dec} & \textbf{Spec Vectors} & \textbf{Throughput} \\
\midrule
PRESENT-80, CPU sequential & 64 / 80  & Round-trip OK & 4/4 Pass & $\approx$0.64~MB/s \\
PRESENT-80, GPU RTX~4070M  & 64 / 80  & Round-trip OK & 4/4 Pass & 1.07~GB/s enc, 1.43~GB/s dec \\
SPECK-64, CPU              & 64 / --- & ---           & ---      & Not started \\
SPECK-64, GPU              & 64 / --- & ---           & ---      & Not started \\
AES-128 baseline           & 128 / 128 & ---          & ---      & Pending \\
\bottomrule
\end{tabular}
\end{table}

PRESENT-80 passes all published vectors. For the same $2^{20}$-block workload the GPU
finishes in 7.85~ms versus 12{,}518~ms on the CPU, a roughly 1595$\times$ speedup.
The CPU implementation is intentionally unoptimized and sits at $\approx$4300 cycles/byte.

\section{Challenges and Issues Encountered}
\begin{itemize}
  \item \textbf{Key-schedule bit ordering:} Loading the 80-bit key into
        \texttt{std::bitset<80>} requires careful mapping because key byte 0 is the most
        significant byte. An off-by-one in the bit indexing produces wrong round keys that
        still pass the round-trip test, so the error only surfaces when running the
        official vectors.

  \item \textbf{Round key extraction:} The spec takes the leftmost 64 bits of the
        80-bit register as the round key. Getting the MSB/LSB alignment right between
        the bitset and the \texttt{uint64\_t} required cross-checking against each
        test vector.

  \item \textbf{P-layer boundary case:} The CPU version handles the $j = 63$ fixed
        point with a ternary inside a 0-to-63 loop. The GPU port splits this into a
        0-to-62 loop with a separate bit-63 check, which is easier to verify.

  \item \textbf{CMake and CUDA:} \texttt{CMAKE\_CUDA\_COMPILER} is not resolved
        automatically in all environments. The GPU build currently requires either the
        \texttt{CUDACXX} environment variable or calling \texttt{nvcc} directly.

  \item \textbf{Scope:} PRESENT-80 uses a 64-bit block with an 80-bit key and 32 rounds.
        For SPECK-64, the key width needs to be confirmed from the SIMON/SPECK appendix
        before implementation starts. Gate-count analysis is out of scope and the PRESENT
        reference value of 1570~GE will only be cited for context.
\end{itemize}

\section{Plans for the Next Two Weeks}
\textbf{Week 1 (May 22--28):} Implement SPECK-64 on CPU and GPU from the NSA spec.
Validate against the published appendix vectors. Write CBC and CTR mode wrappers
that work with both ciphers and verify round-trip correctness for each. \\
\textbf{Week 2 (May 29--June 4):} Add an AES-128 baseline using OpenSSL. Run benchmarks
covering GB/s, cycles/byte, and key-schedule cost at 1~KB, 16~KB, and 64~KB payloads.
Compile the comparison table and finish the report and slides.

\textbf{Planned evaluations:} Full vector sweep for SPECK-64 and both mode wrappers.
Median GB/s over 10 GPU runs. GPU vs.\ CPU speedup. \texttt{.text} size comparison
across PRESENT-80, SPECK-64, and AES-128. A Feistel vs.\ SPN trade-off discussion
for constrained devices.

\textbf{Presentation (June 4):} 10--12 slides on cipher architectures, GPU kernel
design choices, benchmark results, and lightweight-crypto analysis. Slides will be
uploaded to the shared Drive before the session.

\begin{table}[h]
\centering
\small
\begin{tabular}{p{0.22\linewidth}p{0.7\linewidth}}
\toprule
\textbf{Member} & \textbf{Tasks} \\
\midrule
Devansh Gupta    & SPECK-64 CPU/GPU implementation, AES-128 baseline, results plots \\
Priyansh Parikh  & CBC/CTR mode wrappers, benchmark scripts, PRESENT-80 test-vector harness \\
Jeevan N V       & SPECK-64 vector validation, build cleanup, slides \\
\bottomrule
\end{tabular}
\end{table}

\section*{Feasibility and Direction}
PRESENT-80 is done and SPECK-64 is the immediate critical path. One to two days of
implementation and one day of debugging is a reasonable estimate. CBC/CTR wrappers are
straightforward once both ciphers pass vectors. If a custom CUDA AES takes too long,
OpenSSL provides a ready-made reference for the baseline. The final deliverables are
passing test vectors for both ciphers in all modes, a GPU vs.\ CPU comparison table,
and a lightweight-crypto analysis for the presentation.

\vspace{0.3em}
\small\textit{References:} Bogdanov et al., PRESENT (CHES 2007). Beaulieu et al.,
SIMON/SPECK (ePrint 2013/404).

\end{document}
